\section{Introduction}

Public vulnerability databases are operational infrastructure for vulnerability management, dependency risk analysis, and security research \cite{NVDUsers2024,Churakova2026VEX,Mandl2026VulnDetection}. A single CVE can appear in multiple sources, and those sources often expose overlapping structured fields: severity, publication date, references, weakness identifiers, and affected-version ranges. CVE-ID alignment is therefore necessary, but it is not sufficient. Once two records are aligned, downstream systems still need to decide whether their field values are equivalent, merely formatted differently, incomplete, temporally shifted, or factually incompatible.

This problem has become more operationally visible as vulnerability publication volume has increased. NIST's 2026 NVD operations update reports changes to NVD enrichment practice under record CVE growth, including prioritization of actively exploited and high-severity vulnerabilities \cite{NVDOperations2026}. That process reality makes cross-source comparison more than a cleanup exercise: downstream users may encounter aligned records whose enriched fields differ because the sources have different curation workflows, schemas, publication paths, and update timing \cite{NVDProcess2026,NVDStatus2026,GHSAAbout2026,GHSARepo2026}.

This paper studies that post-alignment problem for NVD and GHSA vulnerability records. The alignment step is already done: the normalized corpus contains 100,032 NVD records, 28,785 GHSA records, and 8,066 aligned CVE-ID pairs. The remaining problem is not record matching, but how to handle field values that still disagree after alignment.

The current deterministic baseline already shows that disagreement is uneven across fields. After the 2026-07-15 input-integrity refresh, severity contains 1,749 factual conflicts, affected\_versions contains 651 factual conflicts, published contains 1,897 temporal discrepancies, and references is dominated by incompleteness. The cwe\_ids field is mostly equivalent, with a smaller tail of incomplete and factual-conflict cases. The affected\_versions count had previously been reduced from 1,272 to 652 by rule tightening and then from 652 to 651 by filtering NVD CPE matches explicitly marked \texttt{vulnerable=false}. These numbers are baseline outputs, not gold labels.

The central design choice is to separate discrepancy typing from adjudication: triage before truth. Treating every mismatch as a factual conflict would overstate the problem: a same-day timestamp formatting difference is not the same as incompatible severity values, and a one-sided reference list is not the same as a contradiction. The current pipeline therefore uses five working categories: equivalent, representation\_discrepancy, incomplete, temporal\_discrepancy, and factual\_conflict. Only the last category is passed to evidence-constrained adjudication.

The current manuscript answers three research questions at different maturity levels. RQ1 is descriptive: what distribution of discrepancy types appears across aligned NVD-GHSA fields? RQ2 is currently diagnostic: which deterministic typing rules explain the observed binary differences, and what boundaries appear in complete-impact non-human audits of references and CWE candidates, pending a human discrepancy-typing gold set? RQ3 asks what evidence availability and source-support patterns appear for sampled severity and affected\_versions conflict candidates. Severity remains a silver-label prototype; affected\_versions additionally has two prediction-sealed, CVE-disjoint holdouts reviewed by blind Codex agents, but still no human gold.

The contributions are correspondingly scoped. The paper defines the post-alignment structured-field discrepancy task, operationalizes a five-category field taxonomy for NVD-GHSA pairs, reports reproducible baseline distributions and rule-trigger diagnostics, and specifies an abstention-capable evidence scaffold for sampled factual conflicts. The first affected\_versions holdout exposes a protocol failure: evaluating source selection over all agreed rows mixes discrepancy typing with adjudication conditional on a true factual conflict. A second untouched holdout preregisters these as separate endpoints; its highly selective type candidate and branch-based source method do not exceed full-coverage or fixed-source comparators. The implementation therefore contributes an auditable evaluation protocol and failure analysis, not a validated affected\_versions performance improvement. Human gold remains pending throughout.
